
\documentclass{article}
\usepackage{colortbl}
\usepackage{makecell}
\usepackage{multirow}
\usepackage{supertabular}

\begin{document}

\newcounter{utterance}

\centering \large Interaction Transcript for game `bbh', experiment `snarks', episode 26 with reflexion{-}Qwen3.5{-}27B.
\vspace{24pt}

{ \footnotesize  \setcounter{utterance}{1}
\setlength{\tabcolsep}{0pt}
\begin{supertabular}{c@{$\;$}|p{.15\linewidth}@{}p{.15\linewidth}p{.15\linewidth}p{.15\linewidth}p{.15\linewidth}p{.15\linewidth}}
    \# & \multicolumn{2}{c}{Player} && \multicolumn{2}{c}{Game Master} \\
    \hline

    \theutterance \stepcounter{utterance}  
    & & & \multicolumn{4}{p{0.6\linewidth}}{
        \cellcolor[rgb]{0.9,0.9,0.9}{
            \makecell[{{p{\linewidth}}}]{
                \texttt{\tiny{[P1$\langle$GM]}}
                \texttt{You have attempted this same task before and did not succeed. These are the plans you wrote after those attempts.} \\
\texttt{Trial \#1: New plan: My previous attempt failed because I selected option (B) based on the irony of calling trolls "endangered," but the correct answer was likely (A). I misidentified the sarcastic target; in option (A), the phrase "feed the elephants" is a nonsensical or absurd substitution for the common idiom "feed the trolls," creating a humorous, ironic contrast by treating a mundane or impossible action (feeding elephants) as a standard social interaction, or perhaps the sarcasm lies in the absurdity of the statement itself compared to the expected idiom. Alternatively, I may have misjudged the "endangered" claim as the primary sarcastic element when the sarcasm in (A) relies on the deliberate misuse of the idiom to mock the situation. Next time, I will prioritize identifying which option contains a deliberate, absurd deviation from a common idiom or a statement that is clearly false in a way that mocks the subject, rather than assuming the option with the most obvious "ironic" adjective (like "endangered" for trolls) is the answer. I will specifically look for the option where the literal meaning is so absurd or contradictory to reality that it serves as a clear critique, even if the other option also contains irony. In this specific case, I should have recognized that "feed the elephants" is the sarcastic twist on "feed the trolls," making (A) the intended answer, as the absurdity of the substitution is the core of the humor.} \\
\\ 
\texttt{Trial \#2: New plan: My previous failure stemmed from over{-}analyzing the "endangered" descriptor in option (B) and missing the primary sarcastic mechanism: the absurd substitution of the idiom itself. In the provided examples, sarcasm is consistently generated by stating a falsehood that mocks a situation (e.g., claiming ad hominem attacks help, or comparing league punishments to politics). In the final question, option (A) replaces the standard idiom "feed the trolls" with "feed the elephants." This creates a nonsensical, absurd image that mocks the act of engaging with online provocateurs by treating it as a literal, impossible, or ridiculous task (feeding elephants), whereas option (B) uses the correct idiom with a secondary ironic adjective. Next time, I will prioritize identifying the option that fundamentally breaks the expected idiom or logic to create absurdity, rather than looking for subtle ironic adjectives within a correct idiom. I will specifically check if one option swaps a key noun in a common phrase to create a humorous, illogical statement, as this is the stronger indicator of the intended sarcasm in this specific dataset.} \\
\\ 
\texttt{Trial \#3: New plan: My previous attempts failed because I over{-}analyzed the secondary irony (the "endangered" adjective) in the option containing the correct idiom, while missing the primary sarcastic mechanism: the absurd substitution of the idiom itself. In the provided examples, sarcasm is consistently defined by a statement that is a clear falsehood or absurdity used to mock a situation (e.g., claiming ad hominem attacks help, or calling researching interests a "waste"). In the final question, Option (A) replaces the standard idiom "feed the trolls" with "feed the elephants." This creates a nonsensical, absurd image that mocks the act of engaging with online provocateurs by treating it as a literal, impossible, or ridiculous task. Option (B) uses the correct idiom, making the "endangered" claim a standard ironic remark, but Option (A)'s fundamental break from the expected phrase creates a stronger, more distinct layer of absurdity that fits the pattern of the previous examples where the "falsehood" is the core of the sarcasm. Next time, I will prioritize identifying the option that fundamentally breaks the expected idiom or logic to create absurdity, rather than looking for subtle ironic adjectives within a correct idiom. I will specifically check if one option swaps a key noun in a common phrase to create a humorous, illogical statement, as this is the stronger indicator of the intended sarcasm in this specific dataset. I will select Option (A).} \\
\\ 
\\ 
\texttt{Determine which of two sentences is sarcastic.} \\
\\ 
\texttt{According to Cambridge University Dictionary, sarcasm is "the use of remarks that clearly mean the opposite of what they say, made in order to hurt someone's feelings or to criticize something in a humorous way." Sarcastic sentences often contain satirical or ironic utterances, hyperboles, ambivalent or witty remarks.} \\
\\ 
\texttt{Q: Which statement is sarcastic?} \\
\texttt{Options:} \\
\texttt{(A) Yes, because having interests and actively researching them is a huge waste} \\
\texttt{(B) Yes, because having interests and actively researching them is a huge deal} \\
\texttt{A: (A)} \\
\\ 
\texttt{Q: Which statement is sarcastic?} \\
\texttt{Options:} \\
\texttt{(A) No one is going to disagree with you on this. Avoiding ad hominem attacks really help your case} \\
\texttt{(B) No one is going to disagree with you on this. Ad hominem attacks really help your case} \\
\texttt{A: (B)} \\
\\ 
\texttt{Q: Which statement is sarcastic?} \\
\texttt{Options:} \\
\texttt{(A) Consistency in the league's punishments? What do you think this is supposed to be, politics?} \\
\texttt{(B) Consistency in the league's punishments? What do you think this is supposed to be, moral?} \\
\texttt{A: (A)} \\
\\ 
\texttt{Q: Which statement is sarcastic?} \\
\texttt{Options:} \\
\texttt{(A) You're welcome my dear friend. Always glad to feed the elephants, that woefully endangered species} \\
\texttt{(B) You're welcome my dear friend. Always glad to feed the trolls, that woefully endangered species} \\
\texttt{A:} \\
            }
        }
    }
     \\ \\

    \theutterance \stepcounter{utterance}  
    & \multicolumn{4}{p{0.6\linewidth}}{
        \cellcolor[rgb]{0.561,0.737,0.561}{
            \makecell[{{p{\linewidth}}}]{
                \texttt{\tiny{[P1$\rangle$GM]}}
                \texttt{Based on the pattern established in the previous examples and the analysis from the failed trials, here is the deduction:} \\
\\ 
\texttt{1.  **Analyze the Pattern**:} \\
\texttt{    *   **Example 1**: The sarcasm lies in calling a positive action (researching interests) a "huge waste." The statement is a clear falsehood used to mock.} \\
\texttt{    *   **Example 2**: The sarcasm lies in claiming "Ad hominem attacks really help your case." This is a logical falsehood (ad hominem attacks usually hurt a case) used to mock the situation.} \\
\texttt{    *   **Example 3**: The sarcasm lies in comparing league punishments to "politics" (implying inconsistency and corruption), whereas "moral" would be a sincere comparison. The "politics" option creates the satirical critique.} \\
\\ 
\texttt{2.  **Analyze the Target Question**:} \\
\texttt{    *   **Context**: The phrase "feed the trolls" is a common internet idiom meaning to engage with online provocateurs, which usually gives them attention.} \\
\texttt{    *   **Option (B)**: "Always glad to feed the trolls, that woefully endangered species." This uses the correct idiom. The sarcasm here is in the adjective "endangered," implying trolls are rare (which they are not). This is a standard ironic remark.} \\
\texttt{    *   **Option (A)**: "Always glad to feed the elephants, that woefully endangered species." This option replaces the key noun "trolls" with "elephants."} \\
\texttt{        *   "Feeding elephants" is not the idiom for engaging with online haters.} \\
\texttt{        *   By swapping "trolls" with "elephants," the sentence creates a nonsensical, absurd image.} \\
\texttt{        *   The sarcasm here is derived from the **absurdity of the substitution**. It mocks the situation by treating the act of engaging with trolls as if it were a literal, impossible, or ridiculous task (feeding elephants), or it highlights the absurdity of the "endangered" claim by applying it to a completely unrelated animal in a context where it makes no sense.} \\
\\ 
\texttt{3.  **Applying the Lesson from Previous Trials**:} \\
\texttt{    *   The previous trials concluded that the correct answer relies on the **fundamental break from the expected idiom** to create absurdity, rather than just adding an ironic adjective to a correct idiom.} \\
\texttt{    *   In Example 2, the sarcasm was the *falsehood* of the logic (ad hominem helps).} \\
\texttt{    *   In this question, Option (A) creates a stronger, more distinct layer of absurdity by breaking the idiom itself ("feed the elephants" vs "feed the trolls"). This aligns with the pattern where the "falsehood" or "absurdity" is the core mechanism of the sarcasm. Option (B) is a standard ironic comment, but Option (} \\
            }
        }
    }
    & & \\ \\

    \theutterance \stepcounter{utterance}  
    & & & \multicolumn{2}{p{0.3\linewidth}}{
        \cellcolor[rgb]{0.9,0.9,0.9}{
            \makecell[{{p{\linewidth}}}]{
                \texttt{\tiny{[GM$|$GM]}}
                \texttt{Target: (B)} \\
            }
        }
    }
    & & \\ \\

    \theutterance \stepcounter{utterance}  
    & & & \multicolumn{2}{p{0.3\linewidth}}{
        \cellcolor[rgb]{0.9,0.9,0.9}{
            \makecell[{{p{\linewidth}}}]{
                \texttt{\tiny{[GM$|$GM]}}
                \texttt{game\_result = LOSE} \\
            }
        }
    }
    & & \\ \\

    \theutterance \stepcounter{utterance}  
    & & \multicolumn{4}{p{0.6\linewidth}}{
        \cellcolor[rgb]{0.95,0.95,0.95}{
            \makecell[{{p{\linewidth}}}]{
                \texttt{\tiny{[P1$\leftrightarrow$P]}}
                \texttt{Post{-}episode Reflexion (not sent to GM):} \\
\texttt{New plan: My previous attempts failed because I over{-}analyzed the secondary irony (the "endangered" adjective) in the option containing the correct idiom, while missing the primary sarcastic mechanism: the absurd substitution of the idiom itself. In the provided examples, sarcasm is consistently defined by a statement that is a clear falsehood or absurdity used to mock a situation (e.g., claiming ad hominem attacks help, or calling researching interests a "waste"). In the final question, Option (A) replaces the standard idiom "feed the trolls" with "feed the elephants." This creates a nonsensical, absurd image that mocks the act of engaging with online provocateurs by treating it as a literal, impossible, or ridiculous task. Option (B) uses the correct idiom, making the "endangered" claim a standard ironic remark, but Option (A)'s fundamental break from the expected phrase creates a stronger, more distinct layer of absurdity that fits the pattern of the previous examples where the "falsehood" is the core of the sarcasm. Next time, I will prioritize identifying the option that fundamentally breaks the expected idiom or logic to create absurdity, rather than looking for subtle ironic adjectives within a correct idiom. I will specifically check if one option swaps a key noun in a common phrase to create a humorous, illogical statement, as this is the stronger indicator of the intended sarcasm in this specific dataset. I will select Option (A).} \\
            }
        }
    }
    & \\ \\

\end{supertabular}
}

\end{document}
